%% kapa: exam=1 round=36 year=2025 period=2 form=A subject=accounting subject_ko=회계학 seq=26 number=66
%% source: Q-Net 감정평가사 기출문제 (cst00309) artlSeq=5231525 file=raw/1008/5231525/2025년 제36회 감정평가사 1차시험 2교시 문제지 원본.pdf
%% conversion: kordoc 4.24.51 -> tools/gen1.py (auto); review=proofread-sonnet; flags=none
\begin{question}{66}
다음은 20×1년 설립한 (주)감평의 20×2년 법인세 관련 자료이다.
\begin{tcolorbox}[enhanced, colback=white, colframe=black, boxrule=0.4pt, arc=0pt, left=4pt, right=4pt, top=4pt, bottom=4pt]
\begin{tabular}{@{}p{0.72\linewidth}r@{}}
○ 20×2년 법인세비용차감전순이익 & ￦500,000 \\
\multicolumn{2}{@{}l}{○ 20×2년 처음 발생한 세무조정사항} \\
\quad - 접대비한도초과액 & 20,000 \\
\quad - 감가상각비한도초과액 & 100,000 \\
\multicolumn{2}{@{}l}{○ 설립 이후 연도별 법인세율은 20 \%로 일정하다.} \\
\multicolumn{2}{@{}l}{○ 이연법인세자산(부채)의 실현가능성은 거의 확실하다.}
\end{tabular}
\end{tcolorbox}
20×1년 세무조정사항은 정기적금 미수이자 ￦80,000이 있으며 20×2년 소멸한다. (주)감평의 20×2년 포괄손익계산서에 인식할 법인세비용은?
\begin{choices}
  \item ￦100,000
  \item ￦104,000
  \item ￦120,000
  \item ￦124,000
  \item ￦140,000
\end{choices}
\end{question}
\answer{2}
